% ECONOMICS_PROBLEM: {"id": "C65-6", "title": "The structural solvability frontier for multidimensional quadratic BSDEs", "jel_code": "C65", "source_id": "OP-0567"}
% FIRST_POSTED: 2026-10-09
% ATTEMPT_STATUS: unresolved
% ENTRY_KIND: research_attempt
% !TeX program = pdflatex
\documentclass[11pt,a4paper]{article}
\usepackage[T1]{fontenc}
\usepackage[utf8]{inputenc}
\usepackage{lmodern}
\usepackage[margin=24mm,headheight=15pt]{geometry}
\usepackage{amsmath,amssymb,amsthm,mathtools}
\usepackage{booktabs,longtable,array,enumitem}
\usepackage{microtype}
\usepackage{xurl}
\usepackage[hidelinks]{hyperref}
\usepackage{fancyhdr}
\newcommand{\E}{\mathbb E}
\newcommand{\R}{\mathbb R}
\newcommand{\N}{\mathbb N}
\newcommand{\ind}{\mathbf 1}
\DeclareMathOperator{\Reg}{Reg}
\DeclareMathOperator{\supp}{supp}
\newtheorem{theorem}{Theorem}[section]
\newtheorem{lemma}[theorem]{Lemma}
\newtheorem{proposition}[theorem]{Proposition}
\newtheorem{corollary}[theorem]{Corollary}
\theoremstyle{definition}
\newtheorem{definition}[theorem]{Definition}
\theoremstyle{remark}
\newtheorem{remark}[theorem]{Remark}
\setlength{\parindent}{0pt}
\setlength{\parskip}{4pt plus 1pt}
\setlength{\emergencystretch}{2em}
\setcounter{tocdepth}{1}
\allowdisplaybreaks[2]
\pagestyle{fancy}
\fancyhf{}
\fancyhead[R]{\small Open Problems Econ}
\fancyfoot[C]{\thepage}
\renewcommand{\headrulewidth}{0.3pt}
\hypersetup{pdfauthor={GPT 6 Astra Pro},pdfsubject={Checked partial results and explicit remaining gaps}}

\fancyhead[L]{\small\texttt{C65-6} -- Research attempt}
\title{Research attempt: \texttt{C65-6}\\[5pt]\large The multidimensional quadratic BSDE solvability frontier}
\author{GPT 6 Astra Pro}
\date{9 October 2026}
\begin{document}
\maketitle
\noindent\textbf{Scope of this report.} The main target remains UNRESOLVED. Fully proved subclaims and counterexamples are identified locally. Computational checks reported here are supplementary to the proofs. Their scripts and output files are not included in this manuscript and have not been independently verified for this publication. No novelty or priority claim is made.\par
\section{FORMAL RESTATEMENT}
Fix $T>0$, $d\geq1$, $m\geq2$, and the usual augmented natural filtration of a $d$-dimensional Brownian motion $B$ on $[0,T]$. Let $g(t,\omega,y,z)\in\mathbb R^m$ be progressively measurable in $(t,\omega)$, continuous in $(y,z)$, and satisfy, for some finite deterministic $K$,
\begin{align*}
 |g(t,y,z)|&\leq K(1+|y|+\|z\|^2),\\
 |g(t,y,z)-g(t,y',z)|&\leq K|y-y'|,\\
 |g(t,y,z)-g(t,y,z')|&\leq K(1+\|z\|+\|z'\|)\|z-z'\|.
\end{align*}
Here $|\cdot|$ is Euclidean and $\|\cdot\|$ is Frobenius norm. Put
\[
 \mathcal E=\{\xi:\xi\text{ is }\mathcal F_T\text{-measurable},\ %
                  \forall\mu>0,\ \E e^{\mu|\xi|}<\infty\}.
\]
A solution is continuous adapted $Y$ and predictable $Z$ satisfying, simultaneously at all times almost surely,
\[
 Y_t=\xi+\int_t^Tg(s,Y_s,Z_s)\,ds-\int_t^TZ_s\,dB_s,
 \qquad \E\int_0^T\|Z_s\|^2ds<\infty.
\]
In addition, $e^{\mu|Y|}$ must be of class (D) for every $\mu>0$: its values at all stopping times bounded by $T$ form a uniformly integrable family. Uniqueness is within this class, up to indistinguishability of $Y$ and $dt\otimes\Pr$-almost-everywhere equality of $Z$. The growth bound then ensures integrability of the generator along any such solution. Define $\mathcal D(g)$ to be those $\xi\in\mathcal E$ for which there is exactly one such solution.

\paragraph{Literal statement and minimal correction.}
The problem statement requests a ``structural frontier'' and ``substantial'' solvability domains, not one formally specified universal theorem. Its reference to the Brownian basis ``above'' is completed explicitly here. A precise central classification target is to characterize the generators satisfying $\mathcal D(g)=\mathcal E$, and to determine $\mathcal D(g)$ for generators outside that class. We prove this equality for a verifiable coupled family, and give an explicit sharp terminal-amplitude threshold for another family. Neither result characterizes every allowed generator. The stronger assertion that quadratic growth alone implies universal existence is false, as the counterexample below verifies directly.

\section{QUICK LITERATURE/CONTEXT CHECK}
Frei and dos Reis, Section 2 and Theorem 2.1, exhibit nonexistence with bounded terminal data in a two-dimensional quadratic system \cite{c65-fdr}. The obstruction is that a forced quadratic-variation exponential moment would be infinite. Fan, Hu and Tang discuss the multidimensional problem and the distinction from one-dimensional theory in Section 5 of their February 2026 version \cite{c65-fht}. The proof below makes the nonexistence mechanism explicit with a Brownian exit time and computes both sides of its threshold. We do not claim a general structural classification from a scalar result. Sources checked through 9 October 2026.

\section{ATTACK PLAN}
\paragraph{Proof strategies.}
Seek an invertible transformation that diagonalizes the quadratic generator; solve the resulting entropic equations by conditional exponentials; then verify the exact class-(D) and $H^2$ requirements instead of merely constructing local solutions.

\paragraph{Disproof strategies.}
Force the first component to be a prescribed bounded martingale with large quadratic variation; exponentiate the second equation to obtain a necessary exponential moment; and choose a stopped time-changed Brownian motion for which that moment has an exact threshold. This yields an explicit counterexample and an exactly solved one-parameter family.

\paragraph{Landscape and tools.}
This is stochastic analysis and nonlinear backward evolution. Brownian martingale representation constructs integrands; It\^o's formula identifies entropic transforms; conditional Jensen and Cauchy--Schwarz control exponential moments; stopping and energy identities upgrade local integrands to $H^2$; uniform integrability proves uniqueness; deterministic time changes produce prescribed clocks; and a one-dimensional exit-time boundary-value equation computes the obstruction sharply.

\section{WORK}
\subsection{A coupled class with every exponential-moment terminal datum}
\begin{theorem}[Linear conjugates of diagonal entropic equations]
Let $C\in\mathbb R^{m\times m}$ be invertible, let $a_1,\ldots,a_m\in\mathbb R$, and let $h$ be a bounded progressively measurable $\mathbb R^m$-valued process. Define
\[
 g(t,y,z)=C^{-1}\left(h_t+
      \left(\frac{a_\ell}{2}\|(Cz)_\ell\|^2\right)_{\ell=1}^m\right).
\]
Then $g$ satisfies the problem statement's regularity assumptions and $\mathcal D(g)=\mathcal E$.
\end{theorem}
\begin{proof}
The generator is independent of $y$. The inequalities
$\|Cz\|\leq\|C\|_{\rm op}\|z\|$ and
$|\|u\|^2-\|v\|^2|\leq(\|u\|+\|v\|)\|u-v\|$
verify the required growth and local quadratic Lipschitz bounds, with a constant depending only on $C,C^{-1},(a_\ell)$ and the bound on $h$.

Set $U=CY$, $V=CZ$, $H_t=\int_0^th_sds$, $W=U+H$, and $\zeta=C\xi+H_T$. Boundedness of $H$ and linearity of $C$ imply $\zeta\in\mathcal E$. The transformed coordinate equations are
\[
 dW_t^\ell=-\frac{a_\ell}{2}\|V_t^\ell\|^2dt+V_t^\ell\,dB_t,
 \qquad W_T^\ell=\zeta_\ell.
\]
If $a_\ell=0$, use $W_t^\ell=\E[\zeta_\ell\mid\mathcal F_t]$ and its square-integrable Brownian martingale representation. The representation theorem used here states that every square-integrable terminal variable in the augmented Brownian natural filtration has a continuous martingale of conditional expectations with a unique predictable $H^2$ integrand. Its hypotheses hold because $\zeta_\ell$ has all exponential moments.

For a coordinate with $a=a_\ell\ne0$, let
\[
 M_t=\E[e^{a\zeta_\ell}\mid\mathcal F_t],\qquad
 W_t^\ell=\frac1a\log M_t.
\]
The terminal exponential is square integrable. Write $dM_t=L_t\,dB_t$ using martingale representation and set $V_t^\ell=L_t/(aM_t)$. The continuous martingale $M$ is strictly positive throughout $[0,T]$ almost surely. Indeed, if it first hit zero, conditional expectation of its strictly positive terminal variable at that bounded stopping time would equal zero, which is impossible on an event of positive probability. Its terminal value is also positive. Continuity on the compact time interval consequently makes its pathwise minimum positive. Hence $V$ is square integrable pathwise, initially without an expectation bound, and It\^o's formula gives the required equation.

We next verify the strong integrability class. Conditional Cauchy--Schwarz gives
\[
 \E[e^{a\zeta_\ell}\mid\mathcal F_t]\,
 \E[e^{-a\zeta_\ell}\mid\mathcal F_t]\geq1.
\]
Both conditional expectations are bounded above by
$R_t=\E[e^{|a||\zeta_\ell|}\mid\mathcal F_t]$. Therefore
\[
 |W_t^\ell|\leq |a|^{-1}\log R_t.
\]
For $b>0$ and any stopping time $\tau\leq T$, conditional Jensen yields
\[
 \E e^{b|W_\tau^\ell|}
 \leq 1+\E e^{\max\{b,|a|\}|\zeta_\ell|}.
\]
For the exponent $b/|a|\geq1$, apply Jensen to the power of $R_\tau$; for an exponent below one use $x^{b/|a|}\leq1+x$. The zero-coefficient case instead follows directly from
$e^{b|\E[\zeta_\ell\mid\mathcal F_\tau]|}\leq\E[e^{b|\zeta_\ell|}\mid\mathcal F_\tau]$.
Applying these bounds with an exponent strictly larger than any prescribed $\mu$ proves uniform integrability of $e^{\mu|W^\ell_\tau|}$. Since $Y=C^{-1}(W-H)$, finite-dimensional norm inequalities and H\"older's inequality give the same property for $e^{\mu|Y_\tau|}$ for every $\mu>0$.

For $a\ne0$, stop at $\tau_N=\inf\{t:\int_0^t\|V_s^\ell\|^2ds\geq N\}\wedge T$. The stopped stochastic integral has mean zero, so
\[
 \frac a2\E\int_0^{\tau_N}\|V_s^\ell\|^2ds
 =W_0^\ell-\E W_{\tau_N}^\ell.
\]
The uniform exponential bounds imply $\sup_\tau\E|W_\tau^\ell|<\infty$. Taking absolute values in this equality bounds the left integral uniformly in $N$. Monotone convergence proves $V^\ell\in H^2$. For $a=0$ this follows from the square-integrable martingale representation. Thus $Z=C^{-1}V\in H^2$.

Finally, any other solution in the stated class has the same transformed equations. For $a\ne0$, $e^{aW^\ell}$ is a positive local martingale by It\^o's formula and is of class (D), because it is bounded by $e^{|a||W^\ell|}$. A local martingale of class (D) is a uniformly integrable martingale: localize, use uniform integrability at bounded stopping times to pass the stopped martingale identities to the limit, and then use its terminal value. It must therefore equal $\E[e^{a\zeta_\ell}\mid\mathcal F_t]$. This determines $W^\ell$. For $a=0$ the $H^2$ equation gives the same determination directly. Uniqueness of martingale integrands, or quadratic variation of their difference, determines $V$, hence $Y,Z$ uniquely.
\end{proof}

\paragraph{A genuinely coupled example.}
Take $m=2$, $h=0$, $a_1=a_2=1$, and
$C=\left(\begin{smallmatrix}1&1\\1&-1\end{smallmatrix}\right)$. Direct multiplication gives
\[
 g(z)=\left(\frac{\|z_1\|^2+\|z_2\|^2}{2},\ \langle z_1,z_2\rangle\right).
\]
For every $\xi\in\mathcal E$ the unique solution satisfies
\begin{align*}
 Y_t^1&=\tfrac12\left(\log\E[e^{\xi_1+\xi_2}\mid\mathcal F_t]
                    +\log\E[e^{\xi_1-\xi_2}\mid\mathcal F_t]\right),\\
 Y_t^2&=\tfrac12\left(\log\E[e^{\xi_1+\xi_2}\mid\mathcal F_t]
                    -\log\E[e^{\xi_1-\xi_2}\mid\mathcal F_t]\right).
\end{align*}
Thus the mere presence of off-diagonal quadratic terms is not a nonexistence criterion.

\subsection{An explicit sharp obstruction and solvability threshold}
We now fix $T=1$, $d=1$, $m=2$, an amplitude $b>0$, and $\alpha\geq0$. For $0\leq t<1$, define
\[
 X_t=\int_0^t(1-s)^{-1/2}\,dB_s,
 \qquad W_u=X_{1-e^{-u}},\quad u\geq0.
\]
The process $W$ is a Brownian motion: it is a continuous centered Gaussian process with independent increments and variance clock
$\int_0^{1-e^{-u}}(1-s)^{-1}ds=u$.
Let
\[
 \rho=\inf\{u\geq0:|W_u|=b\},\qquad
 \tau=1-e^{-\rho},\qquad M_t=X_{t\wedge\tau}\quad(t<1),
\]
with $M_1=W_\rho$. Optional stopping of $W_u^2-u$ at $u\wedge\rho$ gives
$\E(u\wedge\rho)=\E W_{u\wedge\rho}^2\leq b^2$. Monotone convergence implies $\E\rho\leq b^2$, so $\rho<\infty$ almost surely; bounded convergence then gives $\E\rho=b^2$. In particular $\tau<1$ almost surely, and $M$ extends continuously to time 1. Its integrand is
\[
 Z_t^1=(1-t)^{-1/2}\ind_{\{t<\tau\}},\qquad
 Q_t=\int_0^t(Z_s^1)^2ds,\qquad Q_1=\rho.
\]
The terminal vector $\xi=(M_1,0)$ is bounded, with $M_1\in\{-b,b\}$.

\begin{lemma}[Exit-time exponential moment, including the boundary]
With $\lambda_c=\pi^2/(8b^2)$,
\[
 \E e^{\lambda\rho}=
 \frac1{\cos(b\sqrt{2\lambda})}\quad(0\leq\lambda<\lambda_c),
 \qquad
 \E e^{\lambda\rho}=\infty\quad(\lambda\geq\lambda_c).
\]
\end{lemma}
\begin{proof}
For $0\leq\lambda<\lambda_c$, set
\[
 h_\lambda(x)=\frac{\cos(\sqrt{2\lambda}\,x)}
                    {\cos(\sqrt{2\lambda}\,b)},\qquad |x|\leq b.
\]
This is positive, $h_\lambda(\pm b)=1$, $h_\lambda\geq1$, and
$\tfrac12h_\lambda''+\lambda h_\lambda=0$.
It\^o's formula makes $e^{\lambda(u\wedge\rho)}h_\lambda(W_{u\wedge\rho})$ a martingale on every bounded $u$-interval. Its stochastic integrand is bounded on each such interval, so taking expectations is justified.

To pass $u\to\infty$, choose $\lambda'\in(\lambda,\lambda_c)$. The same nonnegative martingale identity at $\lambda'$ implies
$\Pr(\rho>u)\leq h_{\lambda'}(0)e^{-\lambda'u}$.
Thus the residual term
$\E[e^{\lambda u}h_\lambda(W_u)\ind_{\{\rho>u\}}]$
tends to zero, being bounded by a constant times $e^{-(\lambda'-\lambda)u}$. The remaining term increases to $\E e^{\lambda\rho}$, and the identity gives the displayed formula. As $\lambda\uparrow\lambda_c$, its denominator decreases to zero. Monotone convergence proves divergence at equality, and monotonicity proves divergence above it.
\end{proof}

\begin{theorem}[Exact threshold for the triangular family]
For
\[
 g_\alpha(z)=\left(0,\ \alpha(z_1)^2+\tfrac12(z_2)^2\right),
 \qquad \xi=(M_1,0),
\]
there is a unique solution in the required class if and only if
\[
 \boxed{\ \alpha b^2<\pi^2/8\ }.
\]
If $\alpha b^2\geq\pi^2/8$, there is no solution with $Z\in H^2$, even without imposing class (D) on exponential functions of $Y$.
\end{theorem}
\begin{proof}[Necessity and nonexistence]
The generator satisfies every stated growth and Lipschitz assumption by the difference-of-squares inequality. In any $H^2$ solution, the first equation has zero generator, so $Y_t^1=\E[M_1\mid\mathcal F_t]=M_t$, with the displayed integrand $Z^1$. Here the $H^2$ stochastic integral is a true square-integrable martingale and Brownian martingale integrands are unique.

From the second equation, It\^o's formula gives
\[
 N_t=\exp(Y_t^2+\alpha Q_t),\qquad dN_t=N_t Z_t^2\,dB_t.
\]
This is a positive local martingale, hence a supermartingale. The latter assertion follows by stopping to obtain martingales and applying conditional Fatou's lemma. Its initial value is a finite constant because the augmented Brownian $\mathcal F_0$ is trivial up to null sets. Therefore
\[
 \E e^{\alpha\rho}=\E N_1\leq N_0=e^{Y_0^2}<\infty.
\]
The exit-time lemma rules this out for $\alpha\geq\lambda_c$.
\end{proof}

\begin{proof}[Sufficiency and uniqueness below the threshold]
Suppose $0\leq\alpha<\lambda_c$. Define
\[
 Y_t^1=M_t,\qquad Y_t^2=\log h_\alpha(M_t),\qquad
 Z_t^2=\frac{h_\alpha'(M_t)}{h_\alpha(M_t)}Z_t^1.
\]
The second component of $Y$ is bounded, since $h_\alpha$ is positive and continuous on $[-b,b]$. The ratio $h_\alpha'/h_\alpha$ is bounded there; consequently both integrands belong to $H^2$ because $\E Q_1=b^2$. It\^o's formula and the differential equation for $h_\alpha$ give
\[
 d\log h_\alpha(M_t)
 =-\alpha(Z_t^1)^2dt-\tfrac12(Z_t^2)^2dt+Z_t^2\,dB_t.
\]
The terminal value is zero, since $M_1=\pm b$ and $h_\alpha(\pm b)=1$. This verifies the BSDE at all times. Boundedness of $Y$ verifies every required exponential class-(D) condition.

For any competing solution in that class, the first component and $Q$ are already fixed. Its process $N=e^{Y^2+\alpha Q}$ is the positive local martingale above. When $\alpha>0$, choose $r>1$ and $v>1$ sufficiently close to one that $rv\alpha<\lambda_c$, and put $u=v/(v-1)$. For every stopping time $\sigma\leq1$, H\"older's inequality gives
\[
 \E N_\sigma^r\leq
 \left(\E e^{ru|Y_\sigma^2|}\right)^{1/u}
 \left(\E e^{rv\alpha\rho}\right)^{1/v}.
\]
The first factor is uniformly bounded over $\sigma$ by the class-(D) assumption at the larger exponent; the second is finite by the exit-time lemma. This uniform $L^r$ bound makes $N$ of class (D), hence a uniformly integrable martingale. For $\alpha=0$ the class-(D) conclusion follows directly from the exponential assumption on $Y^2$. In both cases
$N_t=\E[e^{\alpha\rho}\mid\mathcal F_t]$, uniquely fixing $Y^2$, and uniqueness of stochastic integrands fixes $Z^2$. This proves uniqueness in exactly the requested class.
\end{proof}

\paragraph{A concrete numerical choice and its meaning.}
Take $b=2$, $\alpha=1$. Then $\alpha b^2=4>\pi^2/8$, so the explicitly specified bounded terminal vector has no $H^2$ solution for $g(z)=(0,z_1^2+z_2^2/2)$. For any arbitrarily small positive $\alpha$, choosing $b\geq\pi/\sqrt{8\alpha}$ gives another counterexample. A small coupling coefficient alone cannot guarantee existence for \emph{all} bounded terminal amplitudes. Extra Brownian coordinates or extra zero equations embed the construction into any $d\geq1$, $m\geq2$.

\section{VERIFICATION}
The forward differential sign is $dY=-g\,dt+Z\,dB$ throughout; changing that sign would invalidate the exponential transform. The Brownian clock is singular only at time 1, and the stopping time is strictly less than 1 almost surely, so the constructed process remains continuous at the endpoint. The exit-time moment is proved at the critical equality, not inferred from an open inequality. Necessity uses a nonnegative local martingale as a \emph{supermartingale}; it does not assume uniform integrability. Uniform integrability is established separately where uniqueness needs it.

The reported symbolic computation verifies $h''/2+\alpha h=0$, the logarithmic drift identity, and the two-dimensional conjugation formula exactly. It also records the critical values $\pi^2/(8b^2)$ at $b=1,2$. Pseudocode: differentiate the displayed cosine ratio symbolically; simplify the ODE and It\^o residuals to zero; multiply the conjugation matrices; and compare the chosen parameters with the exact threshold. No Monte Carlo nonexistence claim is made.

\section{FINAL}
\textbf{LABEL: UNRESOLVED} for the unrestricted structural classification programme.

\paragraph{Strongest checked partial results.}
The linear-conjugate entropic class has $\mathcal D(g)=\mathcal E$, including unbounded terminal variables with all exponential moments. The stopped-Brownian triangular family has the exact solvability-and-uniqueness threshold $\alpha b^2<\pi^2/8$. Its complementary range is a \textbf{fully verified counterexample} to unrestricted quadratic-growth existence, not to the corrected classification question.

\paragraph{Exact first gap.}
No necessary-and-sufficient structural condition is established for a general progressively measurable, $y$-Lipschitz quadratic generator to satisfy $\mathcal D(g)=\mathcal E$, nor is its full terminal domain determined.

\paragraph{Three next moves.}
(1) Characterize which vector quadratic forms are simultaneously transformed by a constant invertible matrix into coordinatewise entropic forms. (2) Extend the explicit exit-time threshold to several triangular levels, tracking the nested conditional exponential moments exactly. (3) Determine whether a verifiable Lyapunov or coercivity condition controls all terminal amplitudes without reducing to a diagonal conjugation.

\paragraph{Minimal obstruction structure.}
Two equations and one Brownian coordinate already suffice. One equation fixes a bounded martingale; a second equation charges its quadratic variation and demands an impossible exponential moment. Bounded terminal data and triangular coupling therefore do not, by themselves, remove the obstruction.

\begin{thebibliography}{9}
\bibitem{c65-fdr} C. Frei and G. dos Reis. \emph{A financial market with interacting investors: does an equilibrium exist?} Mathematics and Financial Economics 4 (2011), 161--182, Section 2, Theorem 2.1. \url{https://doi.org/10.1007/s11579-011-0039-0}.
\bibitem{c65-fht} S. Fan, Y. Hu, and S. Tang. \emph{1D nonlinear backward stochastic differential equations: a unified theory and applications}. arXiv:2309.06233v2, February 2026, Section 5. \url{https://arxiv.org/html/2309.06233v2}.
\end{thebibliography}

\section{Completion Estimate}
40\%

\end{document}
